\documentclass[a4paper,11pt]{article}
\usepackage[T2A]{fontenc}
\usepackage[utf8]{inputenc}
\usepackage[russian]{babel}
\usepackage{amsmath,amssymb,mathtools}
\usepackage{geometry}
\geometry{left=22mm,right=22mm,top=20mm,bottom=21mm}
\usepackage{microtype}
\usepackage{tikz}
\usetikzlibrary{arrows.meta,positioning,shapes.geometric,calc}
\usepackage{booktabs,tabularx,array}
\usepackage{enumitem}
\usepackage{fancyhdr}
\usepackage{hyperref}
\usepackage{xcolor}
\usepackage{parskip}
\usepackage{needspace}
\definecolor{accent}{HTML}{007F82}
\definecolor{darkaccent}{HTML}{175E61}
\definecolor{textgray}{HTML}{393939}
\definecolor{softgray}{HTML}{767676}
\hypersetup{colorlinks=true,urlcolor=accent,linkcolor=accent,pdfauthor={Лёвин Артём Александрович},pdftitle={Чек-лист: Дроби, проценты и нахождение целого по части}}
\setlength{\parindent}{0pt}
\setlength{\parskip}{6pt}
\linespread{1.07}
\setlist[enumerate]{leftmargin=7mm,itemsep=8pt,topsep=7pt}
\pagestyle{fancy}
\fancyhf{}
\fancyhead[L]{\footnotesize\color{softgray}Дроби, части и проценты}
\fancyhead[R]{\footnotesize\color{softgray}Ксения Васильченко}
\fancyfoot[C]{\footnotesize\thepage}
\renewcommand{\headrulewidth}{0.45pt}
\renewcommand{\headrule}{\hbox to\headwidth{\color{accent}\leaders\hrule height \headrulewidth\hfill}}
\setlength{\headheight}{22pt}
\newcommand{\mysection}[1]{%
  \Needspace{9\baselineskip}%
  \vspace{5pt}{\color{darkaccent}\Large\bfseries #1\par}%
  \vspace{3pt}{\color{accent}\hrule height 0.55pt}\vspace{8pt}}
\newcommand{\myexample}[1]{\vspace{5pt}{\color{darkaccent}\textbf{#1}}\par\vspace{1pt}}
\newcommand{\keypoint}[1]{{\color{darkaccent}\textbf{Главная мысль.}}~#1}
\newcommand{\skill}[1]{\noindent$\square$\enspace #1\par}
\begin{document}
\begin{titlepage}
\thispagestyle{empty}
\centering
\vspace*{31mm}
{\color{darkaccent}\large МАТЕМАТИКА\par}
\vspace{20mm}
{\fontsize{25}{31}\selectfont Чек-лист\par}
\vspace{13mm}
{\color{accent}\fontsize{17.28}{23}\selectfont Нахождение дроби от числа\\[2pt]и числа по его дроби\par}
\vspace{19mm}
{\large Проценты и текстовые задачи\par}
\vspace{8mm}
{\normalsize 6 класс\par}
\vfill
{\large Ксения Васильченко\par}
\vspace{9mm}
{5 ноября 2026 г.\par}
\vspace{20mm}
{\footnotesize Лёвин Артём Александрович\par}
{\footnotesize эксклюзивно для Ксении Васильченко\par}
\vspace{18mm}
\begin{tikzpicture}[x=1cm,y=1cm]
\draw[accent,line width=0.9pt] (-5.0,0)
  .. controls (-3.6,0.75) and (-1.7,-0.65) .. (0,0)
  .. controls (1.5,0.67) and (3.3,-0.62) .. (5.0,0);
\end{tikzpicture}
\vspace{1mm}
\end{titlepage}

\mysection{1. Как связаны целое, доля и часть}
В задачах на дроби \emph{всё исходное количество} принимают за единицу: $1$. Часть этого количества можно описать двумя способами: в штуках, страницах, метрах, килограммах или как долю от целого.

\begin{center}
\renewcommand{\arraystretch}{1.45}
\begin{tabularx}{.94\linewidth}{@{}>{\raggedright\arraybackslash}X>{\raggedright\arraybackslash}X@{}}
\toprule
Что известно & Как найти неизвестное\\
\midrule
Целое и доля & \textbf{Часть = целое $\cdot$ доля}\\
Часть и соответствующая ей доля & \textbf{Целое = часть $:$ доля}\\
\bottomrule
\end{tabularx}
\end{center}

\myexample{Пример 1. Найти дробь от числа}
Найдём $\tfrac{3}{5}$ от $40$ страниц:
\[
40\cdot\frac35=\frac{40\cdot3}{5}=24\ \text{страницы}.
\]
\myexample{Пример 2. Восстановить целое}
Известно, что $18$ страниц составляют $\tfrac{3}{20}$ всей книги. Тогда:
\[
18:\frac{3}{20}=18\cdot\frac{20}{3}=120\ \text{страниц}.
\]
\keypoint{Чтобы найти всё число по известной части, делите значение части на её долю. Доля должна быть отлична от нуля.}

\mysection{2. Проценты означают сотые доли}
Один процент --- это одна сотая целого. Поэтому проценты переводят в дробь делением на $100$:
\[
15\%=\frac{15}{100}=\frac3{20},\quad
24\%=\frac{24}{100}=\frac6{25},\quad
30\%=\frac{30}{100}=\frac3{10}.
\]
\myexample{Пример 3. По проценту найти всё количество}
Если $36$ страниц составляют $15\%$ книги, то всего:
\[
36:\frac{15}{100}=36:\frac3{20}=36\cdot\frac{20}{3}=240\ \text{страниц}.
\]
Обязательно уточняйте, \emph{чему именно соответствует указанное число}: проценту от целого, уменьшению или оставшейся части.

\clearpage
\mysection{3. Приводим дроби к удобному виду}
При умножении и делении удобно записывать смешанные и десятичные числа как обыкновенные дроби.
\[
1\frac23=\frac{1\cdot3+2}{3}=\frac53,
\quad 5\frac12=\frac{11}{2},
\quad 0{,}4=\frac4{10}=\frac25.
\]
\myexample{Пример 4. Деление на дробь}
Если $110$ составляет $1\tfrac23$ некоторого числа, то само число равно
\[
110:\frac53=110\cdot\frac35=66.
\]
При делении на ненулевую дробь \emph{умножаем на обратную}:
\[
a:\frac bc=a\cdot\frac cb\qquad (b\ne0,\ c\ne0).
\]

\mysection{4. Переводим слова в дроби}
\begin{center}
\renewcommand{\arraystretch}{1.4}
\begin{tabularx}{.94\linewidth}{@{}>{\raggedright\arraybackslash}X>{\raggedright\arraybackslash}X@{}}
\toprule
Выражение в условии & Математическая запись\\
\midrule
«пятая часть» & $\frac15$\\
«третья часть» & $\frac13$\\
«половина» & $\frac12$\\
«половина и ещё четверть» & $\frac12+\frac14=\frac34$\\
\bottomrule
\end{tabularx}
\end{center}

\myexample{Пример 5. Две части описаны словами}
Ученик прочитал половину книги и ещё четверть книги --- всего $102$ страницы. Значит, прочитанная доля равна
\[
\frac12+\frac14=\frac24+\frac14=\frac34.
\]
В книге было
\[
102:\frac34=102\cdot\frac43=136\ \text{страниц}.
\]
\keypoint{Сначала переведите текст в математическую запись. Только после этого выбирайте действие.}

\mysection{5. Короткая памятка по вычислениям}
\begin{itemize}[leftmargin=6mm,itemsep=6pt]
\item \textbf{Сокращайте до умножения:} $\frac6{13}\cdot19\frac12=\frac6{13}\cdot\frac{39}{2}=9$.
\item \textbf{При делении на десятичную дробь} можно умножать делимое и делитель на одинаковую степень десяти:
\[
13{,}5:0{,}09=\frac{13{,}5}{0{,}09}=\frac{1350}{9}=150.
\]
\item \textbf{Деление на $10$, $100$, $1000$} переносит запятую влево на $1$, $2$, $3$ разряда:
\[
12:10=1{,}2,\quad 12:100=0{,}12,\quad 12:1000=0{,}012.
\]
\item \textbf{Проверяйте смысл ответа:} если известна только часть, составляющая меньше единицы, найденное целое должно быть больше этой части.
\end{itemize}

\clearpage
\mysection{6. Оставшаяся часть: вычитание из единицы}
Если продали $\tfrac{3}{11}$ всей массы картофеля, то осталось
\[
1-\frac3{11}=\frac{11}{11}-\frac3{11}=\frac8{11}.
\]
\begin{center}
\begin{tikzpicture}[x=.96cm,y=1cm]
  \draw[darkaccent,line width=1.6pt] (0,0)--(3,0);
  \draw[accent,line width=1.6pt] (3,0)--(11,0);
  \foreach \x in {0,...,11} {\draw[softgray,line width=.45pt] (\x,-.13)--(\x,.13);}
  \node[anchor=south] at (1.5,.18) {$\frac3{11}$ продано};
  \node[anchor=south] at (7,.18) {$\frac8{11}$ осталось};
  \node[anchor=north] at (5.5,-.23) {Вся масса --- $1$};
\end{tikzpicture}
\end{center}
\myexample{Пример 6. Найти общую массу}
После продажи $\tfrac{3}{11}$ картофеля осталось $5\tfrac12$ кг. Оставшаяся доля равна $\tfrac8{11}$, поэтому исходная масса:
\[
5\frac12:\frac8{11}=\frac{11}{2}\cdot\frac{11}{8}
=\frac{121}{16}=7\frac9{16}\ \text{кг}.
\]
\keypoint{В подобных задачах полезно сопоставлять на одной схеме доли и соответствующие им килограммы или метры.}

\mysection{7. Задачи на уменьшение в процентах}
Здесь особенно важно различать \emph{величину уменьшения} и \emph{значение после уменьшения}.

\myexample{Случай А. Известно, на сколько уменьшилась цена}
Во время распродажи товар подешевел на $24\%$. Уменьшение составило $57{,}6$ рубля. Найдём первоначальную цену:
\[
24\%=\frac6{25},\qquad
57{,}6:\frac6{25}=57{,}6\cdot\frac{25}{6}=240\ \text{руб.}
\]
Цена после скидки:
\[
240-57{,}6=182{,}4\ \text{руб.}
\]
Здесь $57{,}6$ рубля соответствует именно \emph{$24\%$ первоначальной цены}.

\myexample{Случай Б. Известно значение после уменьшения}
На новом комбайне работу выполнили за $56$ часов, что на $30\%$ меньше времени, чем на старом. Значит, $56$ часов составляют
\[
100\%-30\%=70\%=\frac7{10}
\]
от прежнего времени. Тогда на старом комбайне требовалось
\[
56:\frac7{10}=56\cdot\frac{10}{7}=80\ \text{часов}.
\]
Проверка: $30\%$ от $80$ часов --- это $24$ часа; $80-24=56$ часов.

\keypoint{Если известно само уменьшение, оно соответствует проценту уменьшения. Если известно новое значение, оно соответствует оставшимся процентам.}

\clearpage
\thispagestyle{fancy}
\begin{center}
{\color{darkaccent}\Large\bfseries Алгоритм: как работать с целым и частью}\par
\vspace{11mm}
\end{center}
\begin{center}
\begin{tikzpicture}[
  >=Latex,
  action/.style={draw=accent,rounded corners=2mm,align=center,text width=9.9cm,minimum height=1.2cm,inner sep=5pt},
  decision/.style={diamond,draw=accent,aspect=2.5,align=center,text width=3.1cm,inner sep=3pt},
  branch/.style={draw=darkaccent,rounded corners=2mm,align=center,text width=5.1cm,minimum height=1.4cm,inner sep=4pt},
  terminal/.style={ellipse,draw=darkaccent,align=center,minimum width=6cm,minimum height=.8cm},
  arr/.style={->,thick,draw=textgray},
  every node/.style={font=\normalsize}
]
\node[terminal] (start) at (0,0) {Прочитайте условие};
\node[action] (a) at (0,-2.0) {Примите целое за $1$. Переведите проценты и слова в дроби};
\node[action] (b) at (0,-4.25) {Найдите долю именно той части, о которой известно число};
\node[decision] (c) at (0,-7.0) {Известно целое?};
\node[branch] (l) at (-3.7,-9.6) {Найдите часть:\\целое $\cdot$ доля};
\node[branch] (r) at (3.7,-9.6) {Найдите целое:\\часть $:$ доля};
\node[action] (check) at (0,-12.1) {Проверьте долю, единицы измерения и смысл ответа};
\node[terminal] (end) at (0,-14) {Запишите ответ};
\draw[arr] (start) -- (a);
\draw[arr] (a) -- (b);
\draw[arr] (b) -- (c);
\draw[arr] (c.west) -| node[pos=.33,above] {да} (l.north);
\draw[arr] (c.east) -| node[pos=.33,above] {нет} (r.north);
\draw[arr] (l.south) -- ++(0,-.45) -| (check.west);
\draw[arr] (r.south) -- ++(0,-.45) -| (check.east);
\draw[arr] (check) -- (end);
\end{tikzpicture}
\end{center}

\clearpage
\mysection{Тренировочный блок \hfill Путь А}
\textcolor{softgray}{10 заданий. Решайте последовательно: от простого к более сложному. В текстовых задачах сначала обозначьте целое и соответствующую долю.}
\vspace{3mm}
\begin{enumerate}[label=\textbf{\arabic*.},leftmargin=8mm,itemsep=14pt]
\item Найдите $\frac38$ от числа $64$.
\item $18$ страниц составляют $\frac37$ книги. Сколько всего страниц в книге?
\item $1\frac14$ некоторого числа равно $45$. Найдите это число.
\item Найдите $12\%$ от числа $250$.
\item $36$ страниц составляют $15\%$ книги. Сколько страниц в книге?
\item Ученик прочитал пятую часть книги --- $16$ страниц. Сколько страниц в книге?
\item Ученик прочитал половину книги и ещё четверть книги. Всего он прочитал $45$ страниц. Сколько страниц в книге?
\item Столб заглублён в землю на $\frac29$ своей длины. Над землёй находится $5\frac14$ м столба. Найдите полную длину столба. Выполните схему.
\item Во время распродажи цена товара уменьшилась на $18\%$, причём уменьшение составило $54$ рубля. Найдите первоначальную цену и цену после скидки.
\item Новый комбайн выполняет работу за $63$ часа, что на $30\%$ меньше времени, чем требовалось старому. Сколько часов требовалось старому комбайну?
\end{enumerate}

\clearpage
\mysection{Ответы к тренировочному блоку}
\vspace{5mm}
\renewcommand{\arraystretch}{1.6}
\begin{center}
\begin{tabularx}{.94\linewidth}{@{}>{\centering\arraybackslash}p{17mm}X@{}}
\toprule
№ & Ответ\\
\midrule
1 & $24$\\
2 & $42$ страницы\\
3 & $36$\\
4 & $30$\\
5 & $240$ страниц\\
6 & $80$ страниц\\
7 & $60$ страниц\\
8 & $6\frac34$ м\\
9 & $300$ руб. первоначально; $246$ руб. после скидки\\
10 & $90$ часов\\
\bottomrule
\end{tabularx}
\end{center}
\vspace{14mm}
{\color{darkaccent}\large Итоговая самопроверка}\par
\vspace{5mm}
\skill{Могу найти дробь от числа умножением.}
\skill{Могу найти целое по известной части делением.}
\skill{Умею переводить проценты и словесные доли в дроби.}
\skill{Вычисляю оставшуюся долю как $1$ минус известная доля.}
\skill{Различаю, дано ли уменьшение или новое значение после уменьшения.}
\vfill
\begin{center}
{\footnotesize\color{softgray}Лёвин Артём Александрович --- эксклюзивно для Ксении Васильченко}
\end{center}
\end{document}
